\paragraph{Accrue interest before changing the principal it prices.}
The same company book can now join its operating cycle to borrowing, repayment,
interest, capital contributions and dividends. The illustrative borrowing has a
fixed annual simple rate of 7.3 percent and an explicit actual-days/365 convention.
These are declared synthetic terms. They provide an exact timing example; fees,
effective-interest valuation, compounding, contractual maturity and source
admission require their own methods.

Interest for each declared interval is principal times the annual rate times its
calendar days divided by 365. A date stored in company memory marks the first day
not yet accrued. Each interval starts at that date and ends exclusively at its
declared boundary. The event occurs at the end of the preceding calendar day and
advances the stored date to the boundary. Principal changes occur at the beginning
of their event date and require that same date in memory. Thus a repayment cannot
change the principal used to price an earlier, unaccrued interval. A second memory
value must equal signed ledger debt; an unexplained principal posting prevents the
next financing method.

Issue USD~100 on 1 January 2027. The 90 days to 1 April accrue
$100\times0.073\times90/365=1.8$ dollars of interest. At that boundary, repay
USD~40 and borrow USD~20: principal becomes USD~80. The next 91 days accrue
USD~1.456, and the remaining 184 days accrue USD~2.944. Total interest is USD~6.2.
The stored date reaches 1 January 2028, covering all 365 days. Repeating or skipping
an interval, changing its event date, or changing principal before the preceding
days have accrued is refused. The convention counts a leap day as another day
while retaining the stated denominator of 365.

The method builder calculates each fraction from the exact dates and embeds it
in the typed definition. The caller supplies neither an interest amount nor an
unchecked year fraction. Rate and elapsed years have inverse dimensions, leaving
a currency amount when multiplied by principal. Each interval remains one finite
declared method; this is not yet a general instrument calendar. A report or
provisional carry does not itself accrue omitted interest. Their balanced totals
therefore remain separate from whether the stored date covers the reporting cutoff.

\paragraph{Keep financing, profit and direct equity distinct.}
Debt receipts and repayments change cash and principal. Interest accrual changes
expense and interest payable; payment reduces that payable and cash. A capital
contribution uses supplied shares times price. Dividend declaration reduces
retained equity and creates a payable; payment settles it. The contribution and
declaration need explicit explanations in the equity statement. A supplied
applicability finding remains unresolved when unknown and unsupported when false.
It does not establish a legal right to distribute retained earnings. Payments
cannot exceed the recorded obligation or available cash.

Combine these financing events with the preceding six operating events. Add a
USD~30 contribution, declare USD~12 of dividends, pay USD~5 of those dividends
and USD~4 of interest. Cash closes at USD~208, receivables at USD~15 and inventory
at USD~25. Debt is USD~80, unpaid interest USD~2.2 and unpaid dividends USD~7,
alongside supplier and payroll payables of USD~12. Profit is USD~8.8 and equity
USD~146.8: assets of USD~248 equal liabilities of USD~101.2 plus that equity.
With interest payments explicitly classified as operating, operating cash is
USD~3 and financing cash USD~105. Selecting the supported financing classification
for interest changes the category, while preserving total cash and the expense.

Every monetary transition is compared with the retained corporate engine under
the same premises. Separate calculations check the amounts and date coverage;
the older engine's acceptance alone cannot establish timing because it accepts a
supplied year fraction. Two charts and matched text produce the same financial
results. Unpaid interest and dividends remain visible across provisional carry;
maturity, valuation, outstanding-share reconciliation, lawful distributions and
complete disclosures remain unresolved.

\paragraph{Compose bounded definitions without weakening each method's checks.}
The combined registry exceeded the original record's structural limit before any
event ran. The successor \texttt{CompanyMethodRegistry/2} checks each method with
the unchanged earlier materializer, including unused methods. Its enclosing
record retains the one-megabyte input limit and the 32-entry limit. Each method
retains its own depth, scalar and 4,096-node bounds. The complete registry identity
remains bound to every call; selecting a valid method cannot conceal a malformed
unused definition.

\texttt{DeclaredCompanyProgram/2} applies the same component-by-component validation
to the opening, schemas, registry, statement map and commands. Its 64-command and
8\,MiB trace bounds remain. \texttt{DeclaredCompanyText/2} uses the same command
grammar and exact hashed imports with the successor registry and runner.
Compatible smaller programs retain the complete monetary transition and report
results; only the declared composition version and its provenance change.
This permits the demonstrated common book while leaving scalability and durable
execution for separate validation.

